
\documentclass{article}
\usepackage{colortbl}
\usepackage{makecell}
\usepackage{multirow}
\usepackage{supertabular}

\begin{document}

\newcounter{utterance}

\centering \large Interaction Transcript for game `wordle', experiment `medium\_frequency\_words\_no\_clue\_no\_critic', episode 6 with Qwen3.6{-}35B{-}A3B{-}FP8{-}without{-}reasoning.
\vspace{24pt}

{ \footnotesize  \setcounter{utterance}{1}
\setlength{\tabcolsep}{0pt}
\begin{supertabular}{c@{$\;$}|p{.15\linewidth}@{}p{.15\linewidth}p{.15\linewidth}p{.15\linewidth}p{.15\linewidth}p{.15\linewidth}}
    \# & \multicolumn{2}{c}{Player} && \multicolumn{2}{c}{Game Master} \\
    \hline

    \theutterance \stepcounter{utterance}  
    & & & \multicolumn{4}{p{0.6\linewidth}}{
        \cellcolor[rgb]{0.9,0.9,0.9}{
            \makecell[{{p{\linewidth}}}]{
                \texttt{\tiny{[P1$\langle$GM]}}
                \texttt{Du bist ein Sprachzauberer, der gerne Wörter anhand der gegebenen Regeln errät.} \\
\\ 
\texttt{Willkommen bei Wordle! Du hast sechs Versuche, um das Zielwort zu erraten, ein gültiges deutsches Wort aus fünf Kleinbuchstaben. Bitte verwende die Tags "erklärung:" und "vermutung:", um für jeden Tipp eine knappe Erklärung zu geben.} \\
\\ 
\texttt{Wenn dein Tipp zum Beispiel "Firma" ist, sollte deine Antwort lauten} \\
\texttt{erklärung: dies ist ein gängiges deutsches Wort mit fünf Buchstaben, und ich beginne meinen Tipp mit diesem Wort.} \\
\texttt{vermutung: Firma} \\
\\ 
\texttt{Nach jedem Tipp wird deine Antwort überprüft, und du erhältst Feedback, das anzeigt, welche Buchstaben korrekt sind (grün), welche Buchstaben korrekt, aber an der falschen Position sind (gelb), und welche Buchstaben falsch sind (rot). Dieses Feedback kann hilfreich sein, um zu bestimmen, welche Buchstaben du in deinem nächsten Tipp einbeziehen oder ausschließen solltest.} \\
\\ 
\texttt{Zum Beispiel könnte das Feedback für "Firma" so aussehen:} \\
\texttt{rückmeldung\_zur\_vermutung: f<gelb> i<gelb> r<grün> m<gelb> a<rot>} \\
\\ 
\texttt{Die Erklärung sollte Details dazu enthalten, wie das rückmeldung\_zur\_vermutung verwendet wird, um zu einem neuen Tipp zu gelangen.} \\
\\ 
\texttt{Erzeuge keinen anderen Text.} \\
\\ 
\texttt{Beginnen wir mit deinem ersten Tipp.} \\
            }
        }
    }
     \\ \\

    \theutterance \stepcounter{utterance}  
    & \multicolumn{4}{p{0.6\linewidth}}{
        \cellcolor[rgb]{0.561,0.737,0.561}{
            \makecell[{{p{\linewidth}}}]{
                \texttt{\tiny{[P1$\rangle$GM]}}
                \texttt{erklärung: Ich beginne mit einem häufigen deutschen Fünf{-}Buchstaben{-}Wort, das eine gute Mischung aus Vokalen und Konsonanten enthält, um schnell Informationen über die Buchstabenverteilung zu erhalten.} \\
\texttt{vermutung: stadt} \\
            }
        }
    }
    & & \\ \\

    \theutterance \stepcounter{utterance}  
    & & & \multicolumn{2}{p{0.3\linewidth}}{
        \cellcolor[rgb]{0.9,0.9,0.9}{
            \makecell[{{p{\linewidth}}}]{
                \texttt{\tiny{[GM$|$GM]}}
                \texttt{attempts = 1} \\
\texttt{target\_word = armee} \\
\texttt{guess = stadt} \\
\texttt{guess\_feedback = s<red> t<red> a<yellow> d<red> t<red>} \\
            }
        }
    }
    & & \\ \\

    \theutterance \stepcounter{utterance}  
    & & & \multicolumn{4}{p{0.6\linewidth}}{
        \cellcolor[rgb]{0.9,0.9,0.9}{
            \makecell[{{p{\linewidth}}}]{
                \texttt{\tiny{[P1$\langle$GM]}}
                \texttt{rückmeldung\_zur\_vermutung: s<rot> t<rot> a<gelb> d<rot> t<rot>} \\
\\ 
\texttt{Gib deine Antwort nur in diesem Format an.} \\
\texttt{erklärung: details} \\
\texttt{vermutung: wort} \\
            }
        }
    }
     \\ \\

    \theutterance \stepcounter{utterance}  
    & \multicolumn{4}{p{0.6\linewidth}}{
        \cellcolor[rgb]{0.561,0.737,0.561}{
            \makecell[{{p{\linewidth}}}]{
                \texttt{\tiny{[P1$\rangle$GM]}}
                \texttt{erklärung: Da 's', 't' und 'd' als falsch markiert sind, schließe ich diese aus. 'a' ist enthalten, aber nicht an Position 3. Ich versuche ein Wort mit 'a' an anderer Stelle und neuen Konsonanten wie 'l', 'm', 'n'.} \\
\texttt{vermutung: lampe} \\
            }
        }
    }
    & & \\ \\

    \theutterance \stepcounter{utterance}  
    & & & \multicolumn{2}{p{0.3\linewidth}}{
        \cellcolor[rgb]{0.9,0.9,0.9}{
            \makecell[{{p{\linewidth}}}]{
                \texttt{\tiny{[GM$|$GM]}}
                \texttt{attempts = 2} \\
\texttt{target\_word = armee} \\
\texttt{guess = lampe} \\
\texttt{guess\_feedback = l<red> a<yellow> m<green> p<red> e<green>} \\
            }
        }
    }
    & & \\ \\

    \theutterance \stepcounter{utterance}  
    & & & \multicolumn{4}{p{0.6\linewidth}}{
        \cellcolor[rgb]{0.9,0.9,0.9}{
            \makecell[{{p{\linewidth}}}]{
                \texttt{\tiny{[P1$\langle$GM]}}
                \texttt{rückmeldung\_zur\_vermutung: l<rot> a<gelb> m<grün> p<rot> e<grün>} \\
\\ 
\texttt{Gib deine Antwort nur in diesem Format an.} \\
\texttt{erklärung: details} \\
\texttt{vermutung: wort} \\
            }
        }
    }
     \\ \\

    \theutterance \stepcounter{utterance}  
    & \multicolumn{4}{p{0.6\linewidth}}{
        \cellcolor[rgb]{0.561,0.737,0.561}{
            \makecell[{{p{\linewidth}}}]{
                \texttt{\tiny{[P1$\rangle$GM]}}
                \texttt{erklärung: 'm' und 'e' sind korrekt an Position 3 und 5. 'a' ist enthalten, aber nicht an Position 2. 'l' und 'p' sind falsch. Ich versuche 'a' an Position 1 oder 4 und nutze neue Buchstaben wie 'n', 'r', 's' (wobei s schon rot war, also besser 'n', 'r', 'k').} \\
\texttt{vermutung: mane} \\
            }
        }
    }
    & & \\ \\

    \theutterance \stepcounter{utterance}  
    & & & \multicolumn{2}{p{0.3\linewidth}}{
        \cellcolor[rgb]{0.9,0.9,0.9}{
            \makecell[{{p{\linewidth}}}]{
                \texttt{\tiny{[GM$|$GM]}}
                \texttt{Error: The length of the guessed word is not 5.} \\
            }
        }
    }
    & & \\ \\

    \theutterance \stepcounter{utterance}  
    & & & \multicolumn{2}{p{0.3\linewidth}}{
        \cellcolor[rgb]{0.9,0.9,0.9}{
            \makecell[{{p{\linewidth}}}]{
                \texttt{\tiny{[GM$|$GM]}}
                \texttt{game\_result = ABORT} \\
            }
        }
    }
    & & \\ \\

\end{supertabular}
}

\end{document}
